\section{Continuous switches and Gaussian insertion}\label{sec:switching}

We work with a microscopic subsequential law supplied by
Proposition~\ref{prop:microscopic-limit}. Its endpoint marks are Gaussian,
whereas genuine graph switches preserve its law. Reversible switches
underlie the random-regular-graph comparison framework of
\citet[Section~3]{BauerschmidtEtAl}. We combine this finite-graph symmetry
with the Gaussian marks to obtain exact orthogonal switching symmetries.
We then use many small orthogonal switches to represent the same
point-process law as a limit of finite matrices with added GOE noise.

\subsection{From permutations to orthogonal switches}

Fix \(k\ge2\), and write \(M=M_k\). For \(O\in O(k)\), define the real matrices
\begin{equation}\label{eq:sw-definitions}
 D_O=\begin{pmatrix}\Id_k&0\\0&O\end{pmatrix},
 \qquad
 B_O=D_OJ_kD_O^\top-J_k.
\end{equation}
In particular, \(B_O\) is symmetric. For an analytic symmetric matrix function
\(M\) with positive definite imaginary part, put
\begin{equation}\label{eq:sw-transformation}
 \mathcal T_O(m,M)=
 \left(
 m-\partial_z\log\det(\Id_{2k}+B_OM),\
 D_O^\top(M^{-1}+B_O)^{-1}D_O
 \right).
\end{equation}
Only the logarithmic derivative is used, so no choice of logarithm is needed.
These expressions are well defined: \(M\) is invertible and
\[
 \Im(M^{-1})=-M^{-*}(\Im M)M^{-1}<0.
\]
Thus \(M^{-1}+B_O\) is invertible, and its inverse again has positive definite
imaginary part.

\begin{proposition}[Continuous switching symmetry]
\label{prop:continuous-switching}
For every fixed \(k\ge2\) and every \(O\in O(k)\), the joint analytic-function
law of \((m,M_k)\) is invariant under \(\mathcal T_O\).
\end{proposition}

\begin{proof}
First suppose that \(O\) is a permutation matrix. In a finite graph, let
\(W=W_{n,k}\) contain the starting endpoint columns followed by the ending
endpoint columns of the \(k\) sampled oriented edges. Except with probability
\(O_k(n^{-1})\), all endpoints are distinct and their induced subgraph is
exactly the marked matching. On this event,
\[
 A\longmapsto A+WB_OW^\top,\qquad W\longmapsto WD_O
\]
permutes the pairing between the two endpoint lists. The resulting graph is
simple and \(d\)-regular. Its marked list has the same restriction, and the
inverse permutation recovers the input. Every graph with its ordered list
has the same joint weight, so this bijection preserves the restricted law.

Write \(R(z)=(A-E-z/n)^{-1}\). The determinant and inverse identities for
finite-rank perturbations give
\begin{align*}
 \widehat m_n(z)
 &=m_n(z)-\partial_z\log\det(\Id_{2k}+B_OM_{n,k}(z)),\\
 \widehat M_{n,k}(z)
 &=D_O^\top(M_{n,k}(z)^{-1}+B_O)^{-1}D_O.
\end{align*}
Indeed \(m_n=n^{-1}\Tr R\), and differentiation in \(z\) supplies exactly
this factor \(n^{-1}\). The formulas also follow without requiring
invertibility of the prelimit compression by using
\((\Id+M_{n,k}B_O)^{-1}M_{n,k}\). In the limit, the positive definite
imaginary part already established makes the inverse formula regular.
Compact-open convergence of analytic functions also gives convergence of
their derivatives. We may therefore pass to the limit, proving the
proposition for permutation matrices.

Next let \(U\in O(k)\), and set \(Q_U=\diag(U,U)\).
The matrices \(C_k\) and \(K_k(m_{\rm r})\) commute with \(Q_U\).
The Gaussian representation~\eqref{eq:marked-representation} consequently
shows that
\[
 \mathcal R_U(m,M)=(m,Q_U^\top M Q_U)
\]
preserves the joint law. This assertion follows even conditionally on
\(\nu\), but we will need switching invariance only unconditionally.
Since \(Q_UD_OQ_U^\top=D_{UOU^\top}\) and
\(Q_UB_OQ_U^\top=B_{UOU^\top}\), the inverse and determinant formulas give
\[
 \mathcal R_U^{-1}\circ\mathcal T_O\circ\mathcal R_U
       =\mathcal T_{UOU^\top}.
\]
Thus the law is invariant under every orthogonal conjugate of a
permutation switch.

Finally, the switching maps form an action:
applying \(\mathcal T_O\) and then \(\mathcal T_{O'}\) gives
\(\mathcal T_{OO'}\). To check this directly, observe that
\begin{equation}\label{eq:sw-composition}
 B_{OO'}=B_O+D_OB_{O'}D_O^\top,\qquad D_{OO'}=D_OD_{O'}.
\end{equation}
Successive inverse updates therefore give the asserted second component.
The determinant ratios for the two updates multiply, which gives the first
component after logarithmic differentiation. A transposition is a
hyperplane reflection. Its orthogonal conjugates are all hyperplane
reflections, and their finite products generate \(O(k)\). The proposition
follows.
\end{proof}

\subsection{The finite Gaussian model}

Fix a discrepancy exponent \(0<\chi<1/4\) from
\eqref{eq:discrepancy}. We use \(r\) planar rotations, each through a
small angle \(\theta\), with aggregate variance of order
\(T=r\theta^2\). The position cutoff \(K\) specifies the finite block
that will receive GOE noise. A much larger temporary cutoff \(L\)
first turns the principal-value resolvent into a finite matrix.
We let \(T\) grow slowly relative to \(K\), and \(K\) slowly relative
to \(r\), so that deleting the outer positions and replacing the
individual summands both have vanishing errors. Choose once and for all
\begin{equation}\label{eq:sw-parameters}
 0<\tau<\frac{\chi}{100000},\qquad
 k=2r,\quad K=r^{1/100},\quad T=r^\tau,\quad
 \theta=\sqrt{T/r},\quad L=r^{100/\chi},
 \qquad r\in\N.
\end{equation}
For each \(r\), take \(O=O_r\) to be the direct sum of \(r\) planar rotations
of angle \(\theta\). Write \(B=B_{O_r}\) and \(M^0=K_k(m_{\rm r})\).
The coefficient below absorbs the constant \(M^0\) in the marked
resolvent when the first-cut argument converts the update into an
additive perturbation.
For all sufficiently large \(r\), the real symmetric matrix
\begin{equation}\label{eq:sw-coefficient}
 \Gamma_r=C_k^{1/2}(B^{-1}+M^0)^{-1}C_k^{1/2}
\end{equation}
is well defined. In fact the singular values of \(B\) are all
\(2|\sin(\theta/2)|\), and \(\|M^0\|\) is bounded independently of \(r\).
Let \(\beta\) be a diagonal matrix of eigenvalues of \(\Gamma_r\), and define
\begin{equation}\label{eq:sw-moments}
 s_1=\Tr\beta,\qquad s_2=\Tr\beta^2.
\end{equation}
These quantities are deterministic.

Index the limiting positions with multiplicity, and put
\[
 I=\{j:|x_j|\le K\},\qquad q_*=|I|,\qquad
 D_I=\diag(x_j:j\in I).
\]
Conditional on \(\nu\), let \(\mathcal W\) be a raw GOE matrix of size \(q_*\):
its off-diagonal entries have variance \(1\), and its diagonal entries have
variance \(2\). It is independent of \(\nu\), apart from its dimension. Define
\begin{equation}\label{eq:sw-inserted-matrix}
 H_r=D_I+s_1\Id_{q_*}+\sqrt{s_2}\,\mathcal W.
\end{equation}
An empty \(I\) gives an empty matrix and zero trace; almost surely this case
does not occur for large \(r\), by \eqref{eq:discrepancy}.

\begin{proposition}[Gaussian insertion]\label{prop:gaussian-insertion}
The coefficients in \eqref{eq:sw-moments} satisfy
\begin{equation}\label{eq:sw-coefficient-bounds}
 \|\beta\|\le C\theta,\qquad
 cT\le s_2\le CT,\qquad |s_1|\le CT,
\end{equation}
with deterministic positive constants. Moreover, as \(r\to\infty\),
\begin{equation}\label{eq:sw-trace-limit}
 m_{\rm r}+\Tr(H_r-z)^{-1}
 \ \Longrightarrow\ m(z)
\end{equation}
jointly at every finite list of points \(z\in\C_+\), and the unscaled
eigenvalue point measure of \(H_r\) converges locally vaguely in law to
\(\nu\). Both convergences are unconditional.

More precisely, generate the Gaussian marks of \(M_k\) conditionally on
\(\nu\), separately for each \(r\), and write
\[
 \mathcal S_r(z)
 =m(z)-\partial_z\log\det(\Id_{2k}+BM_k(z)).
 \]
For almost every \(\nu\), expectations of every bounded smooth function with
bounded derivatives of finitely many real and imaginary trace values differ
by \(o(1)\) when
\(m_{\rm r}+\Tr(H_r-z)^{-1}\) is replaced by \(\mathcal S_r(z)\).
This is a conditional approximation assertion, not conditional switching
invariance.
\end{proposition}

We prove the approximation in three steps. Cutting the principal-value
sums at \(L\) converts the switched trace into the trace of a finite
diagonal matrix perturbed by \(4r\) rank-one terms built from Gaussian
vectors. We then
delete the positions between \(K\) and \(L\), using the signed
principal-value tail cancellation to control the Schur complement.
Finally, on the retained block of order \(K\), a summand-by-summand
replacement turns the perturbation into a scalar shift plus GOE noise
with the same mean and entry covariances.

The graph size \(n\) has already tended to infinity, and
Proposition~\ref{prop:continuous-switching} is exact for every fixed \(k\).
We now let \(r\), and hence \(k=2r\), grow within that limiting law.
No growing-dimensional version of the prelimit eigenvector moment estimates
is needed. For the approximation estimates below, condition on any \(\nu\)
in the probability-one event on which
\eqref{eq:mesoscopic-limit}--\eqref{eq:marked-representation} hold.
Constants may depend on this conditioning input and on a fixed finite list
of spectral parameters. Probability notation in the intermediate estimates
refers to the Gaussian marks.

\subsection{A first cut at a large distance}

We first convert a switched trace into the trace of a finite random matrix.
The elementary consequences of \eqref{eq:discrepancy} that we use are
\begin{equation}\label{eq:sw-tail-sums}
 \nu([-R,R])=2\rho R+O(R^{1-\chi}),\qquad
 \left|\PV\sum_{|x_j|>R}\frac1{x_j}\right|=O(R^{-\chi}),
 \qquad
 \sum_{|x_j|>R}\frac1{x_j^2}=O(R^{-1}).
\end{equation}
The last two follow by partial summation on the positive and negative
half-lines, retaining their constant-density terms together in the
principal-value sum.

Cut the representation of \(M=M_k\) at \(L\), keeping the constant \(M^0\):
\[
 M_L(z)=M^0+
 \sum_{|x_j|\le L}
 \frac{C_k^{1/2}g_jg_j^\top C_k^{1/2}}{x_j-z}.
\]
For fixed \(z\in\C_+\),
\begin{equation}\label{eq:sw-first-tail}
 \|M-M_L\|+\|M'-M_L'\|=O_{\PP}(r^2L^{-\chi}),
 \qquad \|M'\|=O_{\PP}(r^2).
\end{equation}
Here is the dimension dependence in these bounds. The deterministic mean
tail is \(O(L^{-\chi})C_k\) for \(M\), and \(O(L^{-1})C_k\) for its
derivative. Each entry of the centered Gaussian tail has variance
\(O(L^{-1})\), using \eqref{eq:sw-tail-sums}; the derivative has a smaller
variance. There are \((4r)^2\) entries, so the Hilbert--Schmidt estimate gives
\(O_{\PP}(rL^{-1/2})\). This is smaller than the bound displayed in
\eqref{eq:sw-first-tail}. For \(M'\), the absolute scalar weights sum to a
finite constant depending on \(z,\nu\), and the expected sum of the norms of
the positive rank-one matrices is \(O(r)\). The stated \(O_{\PP}(r^2)\) bound is
therefore sufficient and conservative.

We also need a lower bound on the imaginary parts in order to use inverse
updates stably. Keep only the indices \(|x_j|\le r^4\). Their number is
asymptotic to \(2\rho r^4\), and their scalar weights in \(\Im M(z)\) are at
least \(c_z r^{-8}\). If \(b_j=C_k^{1/2}g_j\), Gaussian fourth moments give
\[
 \EE\left\|
 \frac1{|\{j:|x_j|\le r^4\}|}
 \sum_{|x_j|\le r^4}b_jb_j^\top-C_k
 \right\|_{\HS}^2=O(r^{-2}).
\]
Since the smallest eigenvalue of \(C_k\) is bounded below independently of
\(k\), both \(M\) and \(M_L\) consequently satisfy
\begin{equation}\label{eq:sw-imaginary-lower}
 \Im M(z)\ge r^{-5}\Id_{4r},\qquad
 \Im M_L(z)\ge r^{-5}\Id_{4r}
\end{equation}
with probability tending to one. We used \(L\gg r^4\) here.
As \(B^{-1}\) is real symmetric, \eqref{eq:sw-imaginary-lower} bounds the
norms of \((B^{-1}+M)^{-1}\) and \((B^{-1}+M_L)^{-1}\) by \(r^5\).

It follows that cutting both scalar and marked sums in
\[
 \mathcal S_r(z)=m(z)-\Tr(B^{-1}+M(z))^{-1}M'(z)
\]
costs \(o_{\PP}(1)\). To make the margin explicit, the marked tail in
\eqref{eq:sw-first-tail} is \(O_{\PP}(r^{-98})\). The inverse identity bounds
the two trace errors by
\[
 O_{\PP}\bigl(r\cdot r^5\cdot r^{-98}\bigr)
 +O_{\PP}\bigl(r\cdot r^{10}\cdot r^{-98}\cdot r^2\bigr)
 =O_{\PP}(r^{-92}+r^{-85}).
\]
The scalar tail is \(O(L^{-\chi})\).

Let \(D=\diag(x_j:|x_j|\le L)\).
Diagonalizing the real symmetric coefficient \(\Gamma_r\), and rotating
the independent Gaussian row vectors in channel space, gives a matrix
\(\mathcal G\) of independent standard real Gaussian entries with \(4r\)
columns. The finite-rank inverse identity now yields
\begin{equation}\label{eq:sw-finite-matrix}
 \mathcal S_r(z)
 =m_{\rm r}+\Tr(D+A'-z)^{-1}+o_{\PP}(1),
 \qquad A'=\mathcal G\beta\mathcal G^\top .
\end{equation}
Indeed \(B^{-1}+M_L\) is the sum of \(B^{-1}+M^0\) and the marked
compression of \((D-z)^{-1}\), while its derivative is the compression of
\((D-z)^{-2}\). Woodbury's formula gives exactly the trace difference
in \eqref{eq:sw-finite-matrix}.

For completeness, the coefficient estimates
\eqref{eq:sw-coefficient-bounds} follow directly from this construction.
The identity
\[
 (B^{-1}+M^0)^{-1}=(\Id+BM^0)^{-1}B=B+O(\theta^2)
\]
holds in operator norm uniformly in \(r\). Every singular value of this
matrix is comparable to \(\theta\); conjugation by \(C_k^{1/2}\) preserves
this comparison. Since the dimension is \(4r\), we obtain
\(\|\beta\|=O(\theta)\) and \(s_2\asymp r\theta^2=T\). Finally,
\[
 s_1=\Tr(C_kB)+O(r\theta^2),\qquad
 \Tr(C_kB)=2p\,\Tr(O-\Id_k)=O(r\theta^2).
\]
This proves all three bounds, including the cancellation needed for
the mean \(s_1\).

\subsection{Removing the outer positions}

The finite matrix in \eqref{eq:sw-finite-matrix} still has order \(L\).
We next show that its outer positions can be deleted at the level of
resolvent traces. This is necessary before making a Gaussian replacement:
the available small-summand estimate will be applied only to the much
smaller block of order \(K\).

Let \(F=\{j:K<|x_j|\le L\}\), so that the indices of \(D\) split into \(I,F\).
Put
\begin{equation}\label{eq:sw-weights}
 h=T^2,\qquad U_I=(|D_I|+h)^{-1/2},\qquad
 U_F=|D_F|^{-1/2}.
\end{equation}
Partial summation using \eqref{eq:sw-tail-sums} gives
\begin{equation}\label{eq:sw-weight-bounds}
 \Tr U_I^2+\Tr U_F^2=O(\log r),\qquad
 \|U_I^2\|\le h^{-1},\quad \|U_F^2\|\le K^{-1}.
\end{equation}
For example, the inner sum is bounded by integrating the linear count
bound against \((|x|+h)^{-1}\); the far sum is treated in the same way
between \(K\) and \(L\). Both endpoints are fixed powers of \(r\).
In the following estimates, \(\operatorname{polylog}(r)\) denotes a fixed
power of \(\log r\), whose exponent may increase from one occurrence
to the next.

\begin{lemma}[Weighted block estimates]\label{sw:weighted-blocks}
For \(\mathcal A=I,F\), put \(a_I=h\), \(a_F=K\). With probability tending
to one,
\begin{equation}\label{eq:sw-bernstein}
 \left\|U_{\mathcal A}
       (A'_{\mathcal A\mathcal A}-s_1\Id)
       U_{\mathcal A}\right\|
 \le \operatorname{polylog}(r)
       \left(\sqrt{T/a_{\mathcal A}}+\theta\right).
\end{equation}
Moreover, for \(P_0=U_IA'_{IF}U_F\),
\begin{equation}\label{eq:sw-cross-block}
 \|P_0\|_{\HS}^2\le T\operatorname{polylog}(r)
\end{equation}
with probability tending to one.
\end{lemma}

\begin{proof}
The centered matrix in \eqref{eq:sw-bernstein} is a sum of independent
self-adjoint matrices
\[
 \sum_{\alpha=1}^{4r}\beta_{\alpha\alpha}
       (v_\alpha v_\alpha^\top-P),\qquad P=U_{\mathcal A}^2,
\]
where \(v_\alpha\) is Gaussian of covariance \(P\).
Its coefficients may have either sign. Gaussian fourth moments give
\begin{equation}\label{eq:sw-fourth-moment}
 \EE[(v_\alpha v_\alpha^\top)^2]
   =(\Tr P)P+2P^2.
\end{equation}
Truncate each original Gaussian column on the event that all its entries
have absolute value at most \(\log r\), and center the truncated summand.
There are only polynomially many entries, so this alters the matrices with
superpolynomially small probability. Gaussian tails also make the resulting
centering correction superpolynomially small.

Each truncated summand has norm at most
\[
 C\|\beta\|(\log r)^2\Tr P
 \le \theta\operatorname{polylog}(r).
\]
If \(X\) is a self-adjoint random matrix, then
\(\EE[(X-\EE X)^2]=\EE[X^2]-(\EE X)^2\le\EE[X^2]\).
Also truncating \(v_\alpha v_\alpha^\top\) by an event can only decrease
the expectation of its square in positive semidefinite order.
Thus \eqref{eq:sw-fourth-moment}, including for negative
\(\beta_{\alpha\alpha}\), bounds the matrix variance parameter by
\[
 C s_2(\Tr P)\|P\|
 \le \frac{T}{a_{\mathcal A}}\operatorname{polylog}(r).
\]
The bounded self-adjoint matrix Bernstein inequality
\citep[Theorem~1.4]{Tropp} gives, for a sum of independent centered
self-adjoint matrices of norm at most \(b\) and variance parameter
\(\sigma_{\rm mat}^2\), a two-sided tail bounded by
\[
 2|\mathcal A|\exp\left(
 -\frac{c t^2}{\sigma_{\rm mat}^2+bt}\right).
\]
Its dimension factor is polynomial in \(r\). Taking a sufficiently large
power of \(\log r\) in \(t\) proves \eqref{eq:sw-bernstein}.

For the cross block, independence of distinct Gaussian rows and direct
second-moment calculation give
\[
 \EE\|P_0\|_{\HS}^2
 =s_2(\Tr U_I^2)(\Tr U_F^2).
\]
Equation~\eqref{eq:sw-cross-block} follows from
\eqref{eq:sw-weight-bounds} and Markov's inequality, with one additional
logarithmic factor.
\end{proof}

Fix \(z\in\C_+\), and define
\[
 G_F=(D_F+A'_{FF}-z)^{-1},\qquad
 G_I=(D_I+A'_{II}-z)^{-1}.
\]
The next estimates control the influence of the far block even at the fixed
imaginary part \(\Im z\).

\begin{lemma}[Far block and Schur correction]\label{sw:schur}
Let
\[
 \delta_F=\operatorname{polylog}(r)
       \left(\frac{1+T}{K}+\sqrt{\frac TK}+\theta\right).
\]
Then
\begin{equation}\label{eq:sw-far-resolvent}
 \|U_F^{-1}(G_F-D_F^{-1})U_F^{-1}\|\le\delta_F=o(1)
\end{equation}
with probability tending to one, and \(\Tr G_F=o_{\PP}(1)\).
For \(\Sigma=A'_{IF}G_FA'_{FI}\),
\begin{equation}\label{eq:sw-schur-small}
 \|U_I\Sigma U_I\|=o_{\PP}(T^{-5}).
\end{equation}
\end{lemma}

\begin{proof}
Conjugating the matrix to be inverted by \(U_F\) gives
\[
 U_F(D_F+A'_{FF}-z)U_F
 =\operatorname{sign}(D_F)+U_F(A'_{FF}-z)U_F.
\]
The second term has norm at most \(\delta_F\), by
Lemma~\ref{sw:weighted-blocks}, the bound \(|s_1|=O(T)\), and
\eqref{eq:sw-weight-bounds}. Since \(\delta_F=o(1)\), a Neumann expansion
proves \eqref{eq:sw-far-resolvent}. In particular,
\(Q_F=U_F^{-1}G_FU_F^{-1}\) has bounded norm.
By \eqref{eq:sw-tail-sums},
\[
 \left|\sum_{j\in F}\frac1{x_j}\right|=O(K^{-\chi}),
 \qquad
 |\Tr(G_F-D_F^{-1})|
 \le\delta_F\Tr U_F^2=o(1).
\]
This proves the assertion about the far trace.

For the Schur correction, replacing \(G_F\) by \(D_F^{-1}\) costs at most
\begin{equation}\label{eq:sw-schur-replacement}
 \|P_0\|_{\HS}^2\delta_F
 \le T\operatorname{polylog}(r)\delta_F
\end{equation}
in the norm weighted on both sides by \(U_I\).
It remains to estimate
\(A'_{IF}D_F^{-1}A'_{FI}\).
Condition further on the inner Gaussian row block \(\mathcal G_I\), and set
\[
 Q=\mathcal G_I\beta^2\mathcal G_I^\top.
\]
The columns of \(A'_{IF}\), indexed by \(j\in F\), are then independent
centered Gaussian vectors \(a_j\) with covariance \(Q\).
Consequently
\begin{align*}
 \EE\left[\sum_{j\in F}\frac{a_ja_j^\top}{x_j}
       \,\middle|\,\mathcal G_I,\nu\right]
 &=Q\sum_{j\in F}\frac1{x_j},\\
 \Var\left(\sum_{j\in F}\frac{(a_j)_a(a_j)_b}{x_j}
       \,\middle|\,\mathcal G_I,\nu\right)
 &=(Q_{aa}Q_{bb}+Q_{ab}^2)
       \sum_{j\in F}\frac1{x_j^2}.
\end{align*}
Also
\(\EE[Q_{aa}^2\mid\nu]
 =s_2^2+2\Tr\beta^4=O(T^2)\).
Since \(Q\) is positive semidefinite,
\(|Q_{ab}|^2\le Q_{aa}Q_{bb}\); Cauchy--Schwarz bounds the remaining
second moments by \(O(T^2)\).
Multiplying entry \(a,b\) by \((U_I)_{aa}(U_I)_{bb}\), summing its second
moment over \(a,b\), and using \eqref{eq:sw-tail-sums} and
\eqref{eq:sw-weight-bounds}, we obtain
\[
 \EE\left\|
 U_IA'_{IF}D_F^{-1}A'_{FI}U_I
 \right\|_{\HS}^2
 \le C T^2(\log r)^2(K^{-2\chi}+K^{-1}).
\]
Markov's inequality therefore bounds this norm by
\(T\operatorname{polylog}(r)(K^{-\chi}+K^{-1/2})\) with probability
tending to one.

Combining this estimate with \eqref{eq:sw-schur-replacement} gives
\[
 \|U_I\Sigma U_I\|
 \le \operatorname{polylog}(r)\,O\left(
 TK^{-\chi}+TK^{-1/2}+\frac{T(1+T)}K
       +T\sqrt{T/K}+T\theta\right)
\]
with probability tending to one. After multiplication by \(T^5\), all
terms tend to zero by a fixed power of \(r\). For example, the three
potentially slow terms have exponents
\[
 6\tau-\frac{\chi}{100},\qquad
 \frac{13}{2}\tau-\frac1{200},\qquad
 \frac{13}{2}\tau-\frac12,
\]
respectively. They are negative under \eqref{eq:sw-parameters}; the terms
\(T^6K^{-1/2}\) and \(T^6(1+T)/K\) have still ample margins.
This proves \eqref{eq:sw-schur-small}.
\end{proof}

We now convert the small weighted Schur correction into a small trace
error. The fixed-height inner resolvent cannot be estimated directly from
\eqref{eq:sw-bernstein}; instead we first work at the larger height \(h=T^2\).

\begin{lemma}[Inner resolvent and trace cut]\label{sw:trace-cut}
For every fixed \(z\in\C_+\),
\begin{equation}\label{eq:sw-inner-resolvent}
 \|U_I^{-1}G_I(z)U_I^{-1}\|=O_{\PP}(h),
\end{equation}
and
\begin{equation}\label{eq:sw-second-cut}
 \Tr(D+A'-z)^{-1}
 =\Tr(D_I+A'_{II}-z)^{-1}+o_{\PP}(1).
\end{equation}
The same estimates hold jointly for a fixed finite list of \(z\)'s.
\end{lemma}

\begin{proof}
Every diagonal entry of \(U_I(D_I-ih)U_I\) has modulus at least
\(1/\sqrt2\). The added term has norm at most
\[
 \|U_IA'_{II}U_I\|
 \le C T/h+\operatorname{polylog}(r)
       \left(\sqrt{T/h}+\theta\right)=o(1)
\]
with probability tending to one. Hence
\[
 \|U_I^{-1}G_I(ih)U_I^{-1}\|\le C.
\]
The identity \(G_I(ih)^*G_I(ih)=h^{-1}\Im G_I(ih)\) now gives
\[
 \|G_I(ih)U_I^{-1}\|+\|U_I^{-1}G_I(ih)\|
 \le C h^{-1/2}.
\]
Alternatively the same mixed-weight estimates follow from the preceding
weighted inverse and \(\|U_I\|\le h^{-1/2}\).
Use the exact twice-iterated resolvent identity
\begin{align*}
 G_I(z)
 ={}&G_I(ih)+(z-ih)G_I(ih)^2\\
 &+(z-ih)^2G_I(ih)G_I(z)G_I(ih).
\end{align*}
The unweighted middle resolvent in the last term has norm at most
\((\Im z)^{-1}\). Conjugating the identity by \(U_I^{-1}\), the three
terms are bounded by
\[
 C,\qquad C|z-ih|/h,\qquad
 C|z-ih|^2/(h\Im z),
\]
respectively. This proves \eqref{eq:sw-inner-resolvent}.

The Schur complement formula for the full inverse has upper-left block
\(\widetilde G_I=(G_I^{-1}-\Sigma)^{-1}\).
Set \(R_I=U_I^{-1}G_IU_I^{-1}\) and \(S_I=U_I\Sigma U_I\).
Then
\[
 U_I^{-1}\widetilde G_IU_I^{-1}
 =(R_I^{-1}-S_I)^{-1}.
\]
By \eqref{eq:sw-inner-resolvent} and
\eqref{eq:sw-schur-small}, a Neumann expansion gives
\[
 \|U_I^{-1}\widetilde G_IU_I^{-1}\|=O_{\PP}(h),\qquad
 \|U_I^{-1}(\widetilde G_I-G_I)U_I^{-1}\|
 \le O_{\PP}(h^2)o_{\PP}(T^{-5})=o_{\PP}(T^{-1}).
\]
Multiplying the last bound by \(\Tr U_I^2=O(\log r)\) shows that their
traces differ by \(o_{\PP}(1)\).

The full block-inverse trace is
\[
 \Tr\widetilde G_I+\Tr G_F+
 \Tr\bigl(\widetilde G_IA'_{IF}G_F^2A'_{FI}\bigr).
\]
The far trace is \(o_{\PP}(1)\) by Lemma~\ref{sw:schur}.
For the last term recall \(Q_F=U_F^{-1}G_FU_F^{-1}\).
It has bounded norm, and
\[
 U_F^{-1}G_F^2U_F^{-1}=Q_FU_F^2Q_F,\qquad
 \|Q_FU_F^2Q_F\|=O_{\PP}(K^{-1}).
\]
No positivity is asserted for this complex square. The ordinary
operator-norm and Hilbert--Schmidt trace inequality instead gives
\begin{align*}
 \left|\Tr\bigl(\widetilde G_IA'_{IF}G_F^2A'_{FI}\bigr)\right|
 &\le
 \|U_I^{-1}\widetilde G_IU_I^{-1}\|\,
 \|P_0\|_{\HS}^2\,
 \|Q_FU_F^2Q_F\|\\
 &\le \operatorname{polylog}(r)\,O_{\PP}(T^3/K)=o_{\PP}(1).
\end{align*}
This proves \eqref{eq:sw-second-cut}. All estimates used the same Gaussian
matrices, so taking a finite union gives the joint assertion.
\end{proof}

\subsection{Gaussian replacement and convergence of point measures}

We have reduced the switched trace to an inner matrix of size
\[
 q_*=2\rho K+O(K^{1-\chi})\asymp K.
\]
Its perturbation is a sum of \(4r\) independent small rank-one matrices.
Only at this reduced size do we replace that sum by GOE noise.
We use the one-summand replacement and Taylor-expansion method of the
Lindeberg principle; see \citet{ChatterjeeLindeberg} for a general
formulation. The matrix estimate needed here is proved below.

\begin{lemma}[Gaussian replacement]\label{sw:lindeberg}
Conditional on almost every \(\nu\), bounded smooth tests with bounded
derivatives of finitely many trace values have expectation difference
\(o(1)\) between
\[
 D_I+A'_{II}
 \quad\hbox{and}\quad
 H_r=D_I+s_1\Id_{q_*}+\sqrt{s_2}\,\mathcal W.
\]
\end{lemma}

\begin{proof}
Write
\[
 A'_{II}-s_1\Id_{q_*}
 =\sum_{\alpha=1}^{4r}
 \beta_{\alpha\alpha}(g_\alpha g_\alpha^\top-\Id_{q_*}),
\]
where the \(g_\alpha\) are independent standard \(q_*\)-dimensional
Gaussian vectors. The covariance identity
\[
 \Cov(g_ag_b,g_cg_d)
 =\delta_{ac}\delta_{bd}+\delta_{ad}\delta_{bc}
\]
shows that \(gg^\top-\Id\) and a raw GOE matrix have the same entry
covariances. Their means are both zero. Thus replace the summands one at
a time by \(\beta_{\alpha\alpha}\mathcal W_\alpha\), with independent raw
GOE matrices \(\mathcal W_\alpha\). Negative coefficients are allowed:
both covariances are multiplied by \(\beta_{\alpha\alpha}^2\).

For either kind of summand \(X\),
\[
 \EE\|X\|^3\le C\theta^3q_*^3.
\]
For the rank-one term this follows from Gaussian moments of \(\|g\|^2\);
for the GOE term the Hilbert--Schmidt norm gives a sufficient bound.
Fix spectral parameters \(z_1,\ldots,z_b\in\C_+\), and a bounded smooth
test \(\Phi\) on the corresponding real and imaginary trace coordinates,
with bounded derivatives. For any real symmetric matrix \(A\), its
resolvents satisfy \(\|(A-z_\ell)^{-1}\|\le(\Im z_\ell)^{-1}\).
Differentiating the resolvent in a real symmetric direction \(\Delta\)
gives, for \(j=1,2,3\), a trace derivative bounded by
\(C_{z,j}q_*\|\Delta\|^j\). The chain rule therefore bounds the third
directional derivative of the composed test by
\[
 C_{z,\Phi}q_*^3\|\Delta\|^3.
\]
This estimate holds at every intermediate real symmetric matrix,
independently of eigenvalue gaps.

Taylor expansion to second order cancels the means and covariances in
each replacement. Summing the third-order remainders gives total error
\[
 C_{z,\Phi}q_*^6 r\theta^3
 =O\bigl(r^{-0.44+(3/2)\tau}\bigr)=o(1).
\]
Finally \(\sum_\alpha\beta_{\alpha\alpha}\mathcal W_\alpha\) has exactly
the law \(\sqrt{s_2}\,\mathcal W\), proving the lemma.
\end{proof}

\begin{proof}[Completion of Proposition~\ref{prop:gaussian-insertion}]
The coefficient bounds were proved after \eqref{eq:sw-finite-matrix}.
That equation, Lemma~\ref{sw:trace-cut}, and Lemma~\ref{sw:lindeberg} prove
the conditional comparison with \(\mathcal S_r\) in the proposition.
Each comparison concerns bounded tests, so its conditional error can be
averaged by bounded convergence. By
Proposition~\ref{prop:continuous-switching}, the unconditional law of
\(\mathcal S_r\) is the law of \(m\) for every \(r\). This proves
\eqref{eq:sw-trace-limit}; bounded smooth tests suffice to determine
finite-dimensional weak convergence.

We finish by explaining why these trace limits imply convergence of point
measures, without assuming a uniform bound on their total number of
points. Let \(\nu^{(r)}\) be the eigenvalue counting measure of \(H_r\), and
define
\[
 \sigma_r(\dd x)=\frac{\nu^{(r)}(\dd x)}{1+x^2}.
\]
Its total mass is \(\Im\Tr(H_r-i)^{-1}\), whose laws are tight by
\eqref{eq:sw-trace-limit}. Regarding \(\sigma_r\) as a finite positive
measure on the one-point compactification
\(\overline{\R}=\R\cup\{\infty\}\), its laws are therefore tight.
For \(z=e+i\eta\in\C_+\), the kernel
\[
 k_z(x)=(1+x^2)\frac{\eta}{(x-e)^2+\eta^2}
\]
extends continuously to \(\overline{\R}\), with \(k_z(\infty)=\eta\).
Its integral against \(\sigma_r\) is \(\Im\Tr(H_r-z)^{-1}\).
It follows that in any subsequential limit \(\sigma\), the joint law of
these integrals at a countable dense set of \(z\)'s is that of
\(\Im m(z)\). Continuity extends this conclusion to all \(z\).

These integrals determine the de-weighted restriction of \(\sigma\) to
\(\R\). Indeed, for \(f\in C_c(\R)\), integrate the kernel against
\(\pi^{-1}f(e)\dd e\) and let \(\eta\downarrow0\). On \(\R\) the resulting
kernel tends to \((1+x^2)f(x)\), by the Poisson approximate identity.
The integrated kernels are bounded uniformly over
\(x\in\overline{\R}\) and \(0<\eta\le1\): near the compact support this is
the Poisson mass bound, and away from it the factor \(1+x^2\) cancels
the quadratic denominator. At infinity their value is
\(\eta\pi^{-1}\int f(e)\dd e\), which tends to zero.
Dominated convergence is thus applicable to each finite measure
\(\sigma\), including any mass at infinity.

By \eqref{eq:marked-representation}, applying the same inversion to
\(\Im m(z)\) gives \(\int f\,\dd\nu\). Taking a countable determining
family of compactly supported tests, the joint distribution of these
integrals is therefore exactly the law of \(\nu\).
The map from \(\sigma\) to its de-weighted restriction is continuous for
local vague convergence: for a compactly supported \(f\), the function
\((1+x^2)f(x)\), extended by zero at infinity, is continuous on
\(\overline{\R}\). Every subsequential local vague limit of
\(\nu^{(r)}\) consequently has the law of \(\nu\).
This proves the point-process assertion.
\end{proof}
